\section{Justificación del diseño de interacción}

El rediseño del sistema se fundamenta en las heurísticas de Nielsen, las leyes de la Gestalt y los principios de reducción de la carga cognitiva, aplicados de manera complementaria. Las heurísticas orientan la corrección de problemas de usabilidad, mientras que las leyes de la Gestalt justifican la organización visual y jerarquía de la interfaz.

\subsection{Usabilidad — heurísticas de Nielsen}

La propuesta de rediseño del prototipo WebSIS se fundamenta en las heurísticas de usabilidad de \textcite{nielsen1994}. Cada problema identificado en el estudio se relaciona con una solución concreta del prototipo, una heurística aplicable y la evidencia empírica obtenida mediante encuesta, entrevistas y observación directa.

{\scriptsize
\setlength{\tabcolsep}{3pt}
\renewcommand{\arraystretch}{1.25}
\begin{longtable}{|p{3cm}|p{4.2cm}|p{3.4cm}|p{4.7cm}|}
\hline
\textbf{Problema en WebSIS} & \textbf{Solución en el prototipo} & \textbf{Heurística Nielsen} & \textbf{Evidencia del estudio} \\
\hline
\endfirsthead

\hline
\textbf{Problema en WebSIS} & \textbf{Solución en el prototipo} & \textbf{Heurística Nielsen} & \textbf{Evidencia del estudio} \\
\hline
\endhead

P1 -- Sin indicadores de progreso durante la inscripción & Flujo de inscripción en 4 pasos numerados con barra de progreso visible en todo momento & H1 -- Visibilidad del estado del sistema & 76\% encontró la inscripción neutral o difícil. \\
\hline

P2 -- Mensajes de error incomprensibles y sin instrucción correctiva & Mensajes en lenguaje del usuario: qué ocurrió, por qué y qué hacer a continuación & H9 -- Ayuda para reconocer, diagnosticar y recuperarse de errores & E4 perdió una materia por un mensaje ambiguo; 40\% reportó pasos no obvios \\
\hline

P4 -- Menú organizado según lógica técnica interna del DTIC & Separación explícita de roles: Estudiante / Docente / Postulante desde la pantalla principal & H2 -- Correspondencia entre sistema y mundo real & Se realizo una navegacion de 4 seccion en (12 min) \\
\hline

P6 -- Sin acceso directo a tareas frecuentes & Panel de inicio con accesos directos a inscripciones, calificaciones y horarios & H7 -- Eficiencia de uso (aceleradores para usuarios frecuentes) & 96\% usa inscripción; 92\% consulta notas; tareas ocultas a 3 o más clics \\
\hline

P3 -- Interfaz inutilizable en dispositivos móviles & Diseño responsive \textit{mobile-first} con \texttt{width=device-width}; todos los flujos operables desde celular & H7 -- Flexibilidad y eficiencia de uso & E2: ``Desde el celular es imposible''; 24\% usa celular como dispositivo principal \\
\hline

P8 -- Sin soporte contextual dentro del sistema & Sección de avisos institucionales y enlace a mesa de ayuda visibles sin iniciar sesión & H10 - Ayuda y documentación & 80\% necesitó ayuda externa; 68\% acude a un compañero al bloquearse \\
\hline
\caption{Relación entre problemas detectados, soluciones del prototipo y heurísticas de Nielsen}\label{tab:heuristicas-nielsen-prototipo}
\end{longtable}
}
\subsection{Aplicación de leyes de la Gestalt en el prototipo}

Además de las heurísticas de Nielsen, el prototipo incorpora principios perceptivos de la Gestalt para mejorar la organización visual, la lectura rápida y la comprensión de la interfaz. Estas leyes permiten que el estudiante identifique relaciones funcionales entre elementos sin depender únicamente de instrucciones textuales.

{\small
\setlength{\tabcolsep}{4pt}
\renewcommand{\arraystretch}{1.25}
\begin{longtable}{|p{3cm}|p{5.8cm}|p{5.8cm}|}
\hline
\textbf{Ley de Gestalt} & \textbf{Aplicación en el prototipo} & \textbf{Justificación} \\
\hline
\endfirsthead

\hline
\textbf{Ley de Gestalt} & \textbf{Aplicación en el prototipo} & \textbf{Justificación} \\
\hline
\endhead

Proximidad & Agrupación de roles (Estudiante, Docente, Postulante) y bloques del Kárdex (aprobadas, reprobadas, en curso) & Permite identificar rápidamente categorías funcionales sin leer todo el contenido \\
\hline

Continuidad & Flujo de postulación en 4 pasos y orden secuencial de información en Perfil & Guía la mirada de forma natural y reduce la incertidumbre antes de iniciar tareas largas \\
\hline

Figura--Fondo & Formulario de inicio de sesión destacado y mantenimiento de contraste en modo oscuro & Dirige la atención a los elementos accionables y reduce la fatiga visual \\
\hline

Similitud & Estructura visual consistente en comunicados institucionales y módulos principales (Kárdex, Horarios, Perfil) & Facilita el reconocimiento de patrones y acelera el escaneo de información \\
\hline

Región común & Accesos rápidos contenidos en una misma tarjeta en el panel principal & Indica pertenencia funcional y reduce clics para tareas frecuentes \\
\hline

Cierre & Indicadores visuales de progreso académico en el Kárdex & Permite comprender el avance curricular sin requerir lectura detallada \\
\hline
\caption{Aplicación de leyes de la Gestalt en el prototipo WebSIS}\label{tab:leyes-gestalt-prototipo}
\end{longtable}
}
\subsection{Reducción de carga cognitiva}

El rediseño del prototipo también busca reducir la carga cognitiva del estudiante, es decir, la cantidad de esfuerzo mental necesario para comprender la interfaz, recordar rutas, decidir por dónde empezar y verificar si una acción fue completada correctamente. Las siguientes decisiones de diseño transforman problemas detectados en mecanismos concretos de apoyo a la memoria, orientación y confirmación.

{\small
\setlength{\tabcolsep}{4pt}
\renewcommand{\arraystretch}{1.25}
\begin{longtable}{|p{6.7cm}|p{7cm}|}
\hline
\textbf{Situación que incrementa la carga cognitiva} & \textbf{Decisión de diseño en el prototipo} \\
\hline
\endfirsthead

\hline
\textbf{Situación que incrementa la carga cognitiva} & \textbf{Decisión de diseño en el prototipo} \\
\hline
\endhead

El usuario debe memorizar rutas no documentadas cada semestre para navegar el sistema (E3: ``si alguien nuevo lo ve, no entiende nada'') & Acceso directo a las 3 tareas de mayor uso desde la pantalla principal, sin exploración previa \\
\hline

La selección de rol (estudiante / docente / postulante) es ambigua: el sistema no indica por dónde empezar & Etiquetas en lenguaje del usuario con descripción de tarea, por ejemplo: ``Acceso a inscripciones, calificaciones y horarios'' \\
\hline

El flujo de postulación no indica en qué paso está el usuario; el 64\% siente que no recibe información suficiente & Pasos numerados con estado actual visible reducen decisiones activas: el usuario sigue el flujo, no explora \\
\hline
\caption{Decisiones del prototipo orientadas a reducir la carga cognitiva}\label{tab:reduccion-carga-cognitiva}
\end{longtable}
}
